\SourceLeaf{284}%
\[
A\int\frac{\cos y}{m\cos\psi}\mathrm{d}x
\]
taken from \(x=0\) to \(x=h\). We may always set \(\cos y\) equal to one, set \(\cos\psi=1\), and use for \(m\) the value that holds on the line \(M\) or for \(y=0\). The integral thus becomes
\[
\begin{aligned}
&=A\int\frac{\mathrm{d}x}{m^0(1+\mu x+\mu'xx+\text{etc.})}\\
&=\frac{A}{m^0}(h-\tfrac{1}{2}\mu hh+(\tfrac{1}{3}\mu\mu-\tfrac{1}{3}\mu')h^3-\text{etc.})
\end{aligned}
\]
It is always sufficient to use the value that agrees with it through order \(hh\)
\[
\frac{Ah}{\BareRoot{m^0m'}}
\]
in its place.

\Article{13}
The quantities \(l^0,l'\) are determined as follows. Let \(\chi\) be the angle that the great-circle arc \(N\), taken in the direction of increasing \(x\), makes with the meridian at any point of the arc. Take the meridian in the direction from north to south and count the angle from the meridian to the arc from left to right. Let \(S\) be the latitude at that point and \(T\) the longitude reckoned eastward from any chosen meridian. At that point we then have
\[
\begin{aligned}
\mathrm{d}S&=-\cos\chi\mathord{.}\mathrm{d}x+\sin\chi\mathord{.}\mathrm{d}y\\
\mathrm{d}T&=-\frac{\sin\chi}{\cos S}\mathrm{d}x-\frac{\cos\chi}{\cos S}\mathrm{d}y
\end{aligned}
\]
and hence the partial differential quotient
\[
\frac{\mathrm{d}m}{m\mathrm{d}y}=\sin\chi\mathord{.}\frac{\mathrm{d}m}{m\mathrm{d}S}-\frac{\cos\chi}{\cos S}\mathord{.}\frac{\mathrm{d}m}{m\mathrm{d}T}
\]

In our conformal mapping, \(m\) is independent of longitude, or equivalently \(\tfrac{\mathrm{d}m}{m\mathrm{d}T}=0\). Thus
\[
l=\sin\chi\mathord{.}\frac{\mathrm{d}m}{m\mathrm{d}S}
\]

Denote the values of \(\chi\) at \(F\) and \(G\) by \(V^0\) and \(180^0+V'\) (so that in ordinary usage \(V^0\) is the azimuth of the great-circle arc \(FG\) at \(F\), and \(V'\) the azimuth of the great-circle arc
